% problem-000202 5 一氧化氮硫化氢反应
\section*{5. 一氧化氮硫化氢反应}
同为常⻅⼩分⼦，NO星光熠熠， $\mathrm { H } _ { 2 } \mathrm { S }$ 却臭名昭著。然⽽，近期研究发现，在⽣命体系中NO与 $\mathrm { H } _ { 2 } \mathrm { S }$ 窃窃私语。研究⼆者的相互作⽤对于解开⽣命起源、理解⽣理过程的奥秘⾄关重要。NO与 $\mathrm { H } _ { 2 } \mathrm { S }$ 之间的反应及产物⾮常复杂，这⾥我们⼀起关注⼀些重要的基本问题。

\subsection*{5-1}
早在 19 世纪 NO 与 $\mathrm { H } _ { 2 } \mathrm { S }$ 的反应就受到关注。依赖于反应条件甚⾄容器（表⾯可能作为催化剂），可以得到多种产物。最简单的⼀种就是产物中出现两种单质（反应1)，⽽反应过程出现笑⽓也是绕不开的步骤（反应2），条件适当的时候，还会产⽣硫化铵（反应3）。写出反应1-3的⽅程式。

\subsection*{5-2}
研究发现，HNO、HSNO、HSSNO、HS_{2}−等均是⽣理过程的关键物，可由NO与 $\mathrm { H } _ { 2 } \mathrm { S }$ 作⽤产⽣。

\subsection*{5-2}
-1 HSNO(A)表达式给出的原⼦次序就是其连接⽅式。画出A及其共轭碱(B)的Lewis结构式（要求最稳定的形式)。

\subsection*{5-2}
-2 尽管A是最稳定的形式，但是通过低温下光反应，也可以得到其他键合⽅式不同的异构体。画出A所有其他合理异构体的⻣架结构。

\subsection*{5-2}
-3 据信 $\mathrm { S S N O ^ { - } }$ 参与多个⽣理过程。其中涉及SSNO−形成及变化的可能反应及热⼒学常数如下

$
\begin{array}{l l}\mathrm {HSNO + HS_ {2} ^ {-} \rightarrow SSNO^ {-} + HS^ {-} + H^ {+}}&K _ {1} = 5. 0 \times 1 0 ^ {2}\\\mathrm {HSNO + HS_ {2} ^ {-} \rightarrow SSNO + HS^ {-}}&\Delta G _ {2} ^ {\ominus} = - 8. 0 \mathrm{kJmol} ^ {- 1}\end{array}
$

计算 298 K 下 HSSNO 的酸解离常数 $K _ { a } .$ 。

\subsection*{5-3}
为理解⾦属离⼦在相关过程的作⽤，开展了如下研究：空⽓中，络离⼦A（[Ru(EDTA)(OH)]2−, EDTA为⼄⼆胺四⼄酸根）的溶液中，加⼊NaHS溶液并充分搅拌，溶液变为蓝绿⾊，对应于双核络离⼦B的⽣成（反应4）。分别取A和B的溶液，在惰性⽓体保护下，加⼊被NO饱和的磷酸盐缓冲液(pH= 8.2)，充分反应，A转化为C，B转化为D（反应5）。借助于质谱分析，计算出四个络离⼦的式量分别为A: 406.3, B:842.7, C: 419.3, D: 451.4。磁性测量发现，C 为抗磁性，D 为顺磁性。

\subsection*{5-3}
-1 写出络离⼦B、C、D的结构简式。

\subsection*{5-3}
-2 写出反应4和5的离⼦⽅程式。

\subsection*{5-3}
-3 画出C和D中⾦属离⼦�轨道在⼋⾯体场（近似看作正⼋⾯体）中的分裂图并给出�电⼦排布。
